Consider the critical poset of type $P=\qR_3$ (see Table \ref{tab:Loupias frames}) with the following fixed orientation and labelling:
% https://q.uiver.app/#q=WzAsOCxbMCwxLCIxIl0sWzEsMSwiMiJdLFsyLDEsIjMiXSxbMywyLCI1Il0sWzMsMCwiNCJdLFs0LDEsIjYiXSxbNSwxLCI3Il0sWzYsMSwiOCJdLFswLDFdLFsxLDJdLFszLDJdLFsyLDRdLFszLDVdLFs1LDRdLFs2LDVdLFs3LDZdLFszLDQsIiIsMSx7InN0eWxlIjp7ImJvZHkiOnsibmFtZSI6ImRhc2hlZCJ9LCJoZWFkIjp7Im5hbWUiOiJub25lIn19fV1d
\[\begin{tikzcd}[ampersand replacement=\&,sep=tiny]
	\&\&\& 4 \\
	1 \& 2 \& 3 \&\& 6 \& 7 \& 8 \\
	\&\&\& 5
	\arrow[from=2-1, to=2-2]
	\arrow[from=2-2, to=2-3]
	\arrow[from=2-3, to=1-4]
	\arrow[from=2-5, to=1-4]
	\arrow[from=2-6, to=2-5]
	\arrow[from=2-7, to=2-6]
	\arrow[dashed, no head, from=3-4, to=1-4]
	\arrow[from=3-4, to=2-3]
	\arrow[from=3-4, to=2-5]
\end{tikzcd}\]
Its Tits form $q_{\qR_3}$ is
\[
q_{\qR_3}(x) = \sum_{i=1}^{8} x_i^2 -x_1 x_2 - x_2 x_3 - x_3 x_4 - x_5 x_3 -x_5 x_6 - x_6 x_4 - x_7 x_6 -x_8 x_7 + x_5 x_4
\]
because $I_{\qR_3} = \langle \alpha_{34}\alpha_{53} - \alpha_{64}\alpha_{56}\rangle$ so $r_{54}=1$. The associated matrix $B_{\qR_3}$ is given by
\[
B_{\qR_3} = \frac{1}{2}
\small
\begin{pmatrix}
	2 & -1 & 0 & 0 & 0 & 0 & 0 & 0  \\
	-1 & 2 & -1 & 0 & 0 & 0 & 0 & 0  \\
	0 & -1 & 2 & -1 & -1 & 0 & 0 & 0  \\
	0 & 0 & -1 & 2 & 1 & -1 & 0 & 0  \\
	0 & 0 & -1 & 1 & 2 & -1 & 0 & 0  \\
	0 & 0 & 0 & -1 & -1 & 2 & -1 & 0  \\
	0 & 0 & 0 & 0 & 0 & -1 & 2 & -1  \\
	0 & 0 & 0 & 0 & 0 & 0 & -1 & 2  
\end{pmatrix}
\quad \xrightarrow[\text{elimination}]{\text{Gauss-Jordan}} \quad
\begin{pmatrix}
	1 & 0 & 0 & 0 & 0 & 0 & 0 & -1 \\
	0 & 1 & 0 & 0 & 0 & 0 & 0 & -2 \\
	0 & 0 & 1 & 0 & 0 & 0 & 0 & -3 \\
	0 & 0 & 0 & 1 & 0 & 0 & 0 & -2 \\
	0 & 0 & 0 & 0 & 1 & 0 & 0 & -2 \\
	0 & 0 & 0 & 0 & 0 & 1 & 0 & -3 \\
	0 & 0 & 0 & 0 & 0 & 0 & 1 & -2 \\
	0 & 0 & 0 & 0 & 0 & 0 & 0 & 0 
\end{pmatrix}
\]
We have a kernel of dimension $1$ generated by $y_{\qR_3} = (1,2,3,2,2,3,2,1)^t$. Thus, by Lemma \ref{lem:radpossemi} $\Rad(q_{\qR_3})=\ker(B_{\qR_3}) = \langle y_{\qR_3}\rangle$.
We define the representations $B(\lambda)_{\qR_3}$ (with $\lambda\in \K$) as follows:
% https://q.uiver.app/#q=WzAsOCxbMCwxLCJcXEsiXSxbMSwxLCJcXEteMiJdLFsyLDEsIlxcS14zIl0sWzMsMiwiXFxLXjIiXSxbMywwLCJcXEteMiJdLFs0LDEsIlxcS14zIl0sWzUsMSwiXFxLXjIiXSxbNiwxLCJcXEsiXSxbMCwxLCJcXHBoaV97MSwxfSJdLFsxLDIsIlxccGhpX3syLDF9Il0sWzMsMiwiXFxwc2lfezIsMX0iXSxbMiw0LCJcXGJldGFfXFxsYW1iZGEiXSxbMyw1LCJcXHBoaV97MiwxfSIsMl0sWzUsNCwiXFxhbHBoYSIsMl0sWzYsNSwiXFxwc2lfezIsMX0iLDJdLFs3LDYsIlxccHNpX3sxLDF9IiwyXSxbMyw0LCIiLDEseyJzdHlsZSI6eyJib2R5Ijp7Im5hbWUiOiJkYXNoZWQifSwiaGVhZCI6eyJuYW1lIjoibm9uZSJ9fX1dXQ==
\[B(\lambda)_{\qR_3} \defeq 
\begin{tikzcd}[ampersand replacement=\&,sep=small]
	\&\&\& {\K^2} \&\&\& \\
	\K \& {\K^2} \& {\K^3} \&\& {\K^3} \& {\K^2} \& \K \\
	\&\&\& {\K^2}
	\arrow["{\phi_{1,1}}", from=2-1, to=2-2]
	\arrow["{\phi_{2,1}}", from=2-2, to=2-3]
	\arrow["{\beta_\lambda}", from=2-3, to=1-4]
	\arrow["\alpha"', from=2-5, to=1-4]
	\arrow["{\psi_{2,1}}"', from=2-6, to=2-5]
	\arrow["{\psi_{1,1}}"', from=2-7, to=2-6]
	\arrow[dashed, no head, from=3-4, to=1-4]
	\arrow["{\psi_{2,1}}", from=3-4, to=2-3]
	\arrow["{\phi_{2,1}}"', from=3-4, to=2-5]
\end{tikzcd}\]
where $\alpha \defeq \left( \begin{smallmatrix} 1 & 1 & 0 \\ 0 & 1 & 1 \end{smallmatrix} \right)$, $\beta_\lambda \defeq \left( \begin{smallmatrix} \lambda & 1 & 1 \\ 1 & 0 & 1	\end{smallmatrix} \right)$, $\phi_{n,m}\defeq \left( \begin{smallmatrix} \boldsymbol{1}_n \\ \boldsymbol{0}_{m,n}\end{smallmatrix} \right)$ and $\psi_{n,m}\defeq \left( \begin{smallmatrix} \boldsymbol{0}_{m,n} \\ \boldsymbol{1}_n \end{smallmatrix} \right)$ (with $\boldsymbol{1}_n \in \matspace{n}$ identity matrix, $\boldsymbol{0}_{m,n}\in \matspace{m\times n}$ zero matrix).
The relations are preserved since $\beta_\lambda\psi_{2,1} = 
\left( \begin{smallmatrix}
	1 & 1\\
	0 & 1 
\end{smallmatrix}\right)
=\alpha\phi_{2,1}$.

We have $\dimvec{B(\lambda)_{\qR_3}}=y_{\qR_3}$ and we prove the following:

\begin{lemma}\label{lemma:brickcontR3}
	The representations $B(\lambda)_{\qR_3}$ with $\lambda\in \K$ (with $\K$ algebraically closed field) form an infinite family of non-isomorphic bricks with the same dimension. Therefore, the incidence $\K$-algebra $\incidencealgebra{\qR_3}$ is brick-continuous.
\end{lemma}
\begin{proof}
	Let $f\colon B(\lambda)_{\qR_3} \to B(\lambda')_{\qR_3}$ be a morphism of representations.
	% https://q.uiver.app/#q=WzAsMTYsWzAsMSwiXFxLIl0sWzEsMSwiXFxLXjIiXSxbMiwxLCJcXEteMyJdLFszLDIsIlxcS14yIl0sWzQsMCwiXFxLXjIiXSxbNSwxLCJcXEteMyJdLFs2LDEsIlxcS14yIl0sWzcsMSwiXFxLIl0sWzAsNCwiXFxLIl0sWzEsNCwiXFxLXjIiXSxbMiw0LCJcXEteMyJdLFszLDUsIlxcS14yIl0sWzQsMywiXFxLXjIiXSxbNSw0LCJcXEteMyJdLFs2LDQsIlxcS14yIl0sWzcsNCwiXFxLIl0sWzAsMSwiXFxwaGlfezEsMX0iXSxbMSwyLCJcXHBoaV97MiwxfSJdLFszLDIsIlxccHNpX3syLDF9IiwwLHsibGFiZWxfcG9zaXRpb24iOjMwfV0sWzIsNCwiXFxiZXRhX1xcbGFtYmRhIl0sWzMsNSwiXFxwaGlfezIsMX0iLDIseyJsYWJlbF9wb3NpdGlvbiI6NzB9XSxbNSw0LCJcXGFscGhhIiwyXSxbNiw1LCJcXHBzaV97MiwxfSIsMl0sWzcsNiwiXFxwc2lfezEsMX0iLDJdLFszLDQsIiIsMSx7InN0eWxlIjp7ImJvZHkiOnsibmFtZSI6ImRhc2hlZCJ9LCJoZWFkIjp7Im5hbWUiOiJub25lIn19fV0sWzgsOSwiXFxwaGlfezEsMX0iXSxbOSwxMCwiXFxwaGlfezIsMX0iXSxbMTEsMTAsIlxccHNpX3syLDF9Il0sWzEwLDEyLCJcXGJldGFfe1xcbGFtYmRhJ30iLDAseyJsYWJlbF9wb3NpdGlvbiI6NDB9XSxbMTEsMTMsIlxccGhpX3syLDF9IiwyXSxbMTMsMTIsIlxcYWxwaGEiLDJdLFsxNCwxMywiXFxwc2lfezIsMX0iLDJdLFsxNSwxNCwiXFxwc2lfezEsMX0iLDJdLFsxMSwxMiwiIiwxLHsic3R5bGUiOnsiYm9keSI6eyJuYW1lIjoiZGFzaGVkIn0sImhlYWQiOnsibmFtZSI6Im5vbmUifX19XSxbMCw4LCJmXzEiLDIseyJjb2xvdXIiOlswLDEwMCw2MF19LFswLDEwMCw2MCwxXV0sWzEsOSwiZl8yIiwyLHsiY29sb3VyIjpbMCwxMDAsNjBdfSxbMCwxMDAsNjAsMV1dLFsyLDEwLCJmXzMiLDIseyJjb2xvdXIiOlswLDEwMCw2MF19LFswLDEwMCw2MCwxXV0sWzMsMTEsImZfNSIsMCx7ImxhYmVsX3Bvc2l0aW9uIjozMCwiY29sb3VyIjpbMCwxMDAsNjBdfSxbMCwxMDAsNjAsMV1dLFs0LDEyLCJmXzQiLDIseyJsYWJlbF9wb3NpdGlvbiI6NDAsImNvbG91ciI6WzAsMTAwLDYwXX0sWzAsMTAwLDYwLDFdXSxbNywxNSwiZl84IiwwLHsiY29sb3VyIjpbMCwxMDAsNjBdfSxbMCwxMDAsNjAsMV1dLFs2LDE0LCJmXzciLDAseyJjb2xvdXIiOlswLDEwMCw2MF19LFswLDEwMCw2MCwxXV0sWzUsMTMsImZfNiIsMCx7ImNvbG91ciI6WzAsMTAwLDYwXX0sWzAsMTAwLDYwLDFdXV0=
	\[\begin{tikzcd}[ampersand replacement=\&,cramped]
		\&\&\&\& {\K^2} \&\&\& \\
		\K \& {\K^2} \& {\K^3} \&\&\& {\K^3} \& {\K^2} \& \K \\
		\&\&\& {\K^2} \\
		\&\&\&\& {\K^2} \\
		\K \& {\K^2} \& {\K^3} \&\&\& {\K^3} \& {\K^2} \& \K \\
		\&\&\& {\K^2}
		\arrow["{f_4}"'{pos=0.4}, color={rgb,255:red,255;green,51;blue,51}, from=1-5, to=4-5]
		\arrow["{\phi_{1,1}}", from=2-1, to=2-2]
		\arrow["{f_1}"', color={rgb,255:red,255;green,51;blue,51}, from=2-1, to=5-1]
		\arrow["{\phi_{2,1}}", from=2-2, to=2-3]
		\arrow["{f_2}"', color={rgb,255:red,255;green,51;blue,51}, from=2-2, to=5-2]
		\arrow["{\beta_\lambda}", from=2-3, to=1-5]
		\arrow["{f_3}"', color={rgb,255:red,255;green,51;blue,51}, from=2-3, to=5-3]
		\arrow["\alpha"', from=2-6, to=1-5]
		\arrow["{f_6}", color={rgb,255:red,255;green,51;blue,51}, from=2-6, to=5-6]
		\arrow["{\psi_{2,1}}"', from=2-7, to=2-6]
		\arrow["{f_7}", color={rgb,255:red,255;green,51;blue,51}, from=2-7, to=5-7]
		\arrow["{\psi_{1,1}}"', from=2-8, to=2-7]
		\arrow["{f_8}", color={rgb,255:red,255;green,51;blue,51}, from=2-8, to=5-8]
		\arrow[dashed, no head, from=3-4, to=1-5]
		\arrow["{\psi_{2,1}}"{pos=0.3}, from=3-4, to=2-3]
		\arrow["{\phi_{2,1}}"'{pos=0.7}, from=3-4, to=2-6]
		\arrow["{f_5}"{pos=0.3}, color={rgb,255:red,255;green,51;blue,51}, from=3-4, to=6-4]
		\arrow["{\phi_{1,1}}", from=5-1, to=5-2]
		\arrow["{\phi_{2,1}}", from=5-2, to=5-3]
		\arrow["{\beta_{\lambda'}}"{pos=0.4}, from=5-3, to=4-5]
		\arrow["\alpha"', from=5-6, to=4-5]
		\arrow["{\psi_{2,1}}"', from=5-7, to=5-6]
		\arrow["{\psi_{1,1}}"', from=5-8, to=5-7]
		\arrow[dashed, no head, from=6-4, to=4-5]
		\arrow["{\psi_{2,1}}", from=6-4, to=5-3]
		\arrow["{\phi_{2,1}}"', from=6-4, to=5-6]
	\end{tikzcd}\]
	We have that $f_1, f_8$ are just scalar multiplications of elements of $\K$, while $f_2,f_4,f_5,f_7 \in \matspace{2}$ and $f_3,f_6 \in \matspace{3}$. We want to understand the property that $f$ must have, studying the commutativity of the squares in the previous diagram.\smallskip
	
	\noindent $(1)$ The square involving $f_1$ and $f_2$ gives us
	$
	f_2(e_1) =
	\begin{pmatrix}
		f_1 \\
		0
	\end{pmatrix}
	$
	\smallskip
	
	\noindent $(2)$ The square involving $f_2$ and $f_3$ gives us that $\phi_{2,1}f_2(e_i) =f_3(\phi_{2,1}(e_i))$ that means
	\[f_3(e_1) = 
	\begin{pmatrix}
		f_2(e_1) \\
		0
	\end{pmatrix} = \begin{pmatrix}
		f_1 \\
		0 \\
		0
	\end{pmatrix} \qquad 
	f_3(e_2) = 
	\begin{pmatrix}
		f_2(e_2) \\
		0
	\end{pmatrix}\]
	
	\noindent $(3)$ The square involving $f_7$ and $f_8$ gives us
	$
	f_7(e_2) = 
	\begin{pmatrix}
		0 \\
		f_8
	\end{pmatrix}
	$
	\smallskip	
	
	\noindent $(4)$ The square involving $f_6$ and $f_7$ gives us that $\psi_{2,1}f_7(e_i) =f_6(\psi_{2,1}(e_i))$ that means
	\[f_6(e_2) = 
	\begin{pmatrix}
		0 \\
		f_7(e_1)
	\end{pmatrix} \qquad 
	f_6(e_3) = 
	\begin{pmatrix}
		0 \\
		f_7(e_2)
	\end{pmatrix}= \begin{pmatrix}
		0 \\
		0 \\
		f_8
	\end{pmatrix}
	\]
	
	\noindent $(5)$ The square involving $f_3$ and $f_5$ gives us that $\psi_{2,1}f_5(e_i) =f_3(\psi_{2,1}(e_i))$ that means
	\[f_3(e_2) = 
	\begin{pmatrix}
		0 \\
		f_5(e_1)
	\end{pmatrix} = \begin{pmatrix}
	f_2(e_2) \\
	0
	\end{pmatrix} =\begin{pmatrix}
	0 \\ a \\	0
	\end{pmatrix} \qquad 
	f_3(e_3) = 
	\begin{pmatrix}
		0 \\
		f_5(e_2)
	\end{pmatrix}
	\]
	for some $a\in \K$ (we used $(3)$). 
	
	\noindent $(6)$ The square involving $f_5$ and $f_6$ gives us that $\phi_{2,1}f_5(e_i) =f_6(\phi_{2,1}(e_i))$ that means
	\[f_6(e_1) = 
	\begin{pmatrix}
		f_5(e_1) \\ 0
	\end{pmatrix} = \begin{pmatrix}
		a \\ 0 \\ 0
	\end{pmatrix} \qquad 
	f_6(e_2) = 
	\begin{pmatrix}
		f_5(e_2) \\0 
	\end{pmatrix} = \begin{pmatrix}
	0 \\ f_7(e_1) 
	\end{pmatrix} =\begin{pmatrix}
	0 \\ b \\	0
	\end{pmatrix} 
	\]
	for some $b\in \K$ (we used $(4)$).
	
	\noindent $(7)$ The square involving $f_4$ and $f_6$ gives us that $\alpha f_6(e_i) =f_4(\alpha(e_i))$ that means
	\[f_4(e_1) = \alpha f_6(e_1)=
	\begin{pmatrix}
		a \\ 0
	\end{pmatrix} \qquad
	f_4(e_2) = \alpha f_6(e_3)=
	\begin{pmatrix}
		0 \\ f_8
	\end{pmatrix}	 
	\]
	\[\begin{pmatrix}
		a \\ f_8
	\end{pmatrix}=f_4(e_1+e_2) = \alpha f_6(e_2)=
	\begin{pmatrix}
		b \\ b
	\end{pmatrix}\]
	so we conclude that $a=b=f_8$. 
	
	\noindent $(8)$ The square involving $f_3$ and $f_4$ gives us that $\beta_{\lambda'}f_3(e_i) =f_4(\beta_\lambda(e_i))$. For $i=1$ we have that
	\[\begin{pmatrix}
		\lambda f_8 \\ f_8
	\end{pmatrix}=f_4(\lambda e_1+e_2) = \beta_{\lambda'}f_3(e_1)=
	\begin{pmatrix}
		\lambda' f_1 \\ f_1
	\end{pmatrix}\]
	so we conclude that $f_1=f_8$ and $\lambda f_1= \lambda' f_1$. 
	
	Therefore we get 
	\begin{equation}\tag{*}\label{eq:R3maps}
		f_1=f_8 \qquad f_2= f_4=f_5=f_7 =f_1 \boldsymbol{1}_2 \qquad  f_3=f_6= f_1 \boldsymbol{1}_3
	\end{equation}
	
	Now if $\lambda=\lambda'$ we conclude that any morphism $f\colon B(\lambda)_{\qR_3} \to B(\lambda)_{\qR_3}$ depends only on $f_1 \in \K$ by \eqref{eq:R3maps}, therefore $\End(B(\lambda)_{\qR_3})\cong \K$, so  $B(\lambda)_{\qR_3}$ is a brick for any $\lambda\in \K$.
	
	If $\lambda\neq\lambda'$ we have that $\lambda f_1= \lambda' f_1$ and if $f_1\neq 0$ it would mean that $\lambda =\lambda'$, contradiction. Thus, for any $\lambda\neq\lambda'$, we must have $\Hom(B(\lambda)_{\qR_3},B(\lambda')_{\qR_3})=0$, so $B(\lambda)_{\qR_3}$ and $B(\lambda')_{\qR_3}$ cannot be isomorphic.
	
	We conclude that $\{B(\lambda)_{\qR_3}\}_{\lambda\in \K}$ is an infinite family of non-isomorphic bricks with the same dimension and consequently the incidence $\K$-algebra $\incidencealgebra{\qR_3}$ is brick-continuous.
\end{proof}

\Paragraphtitlebf{Construction of $G_{\qR_3}$} We have the brick-continuous family $\{B(\lambda)_{\qR_3}\}_{\lambda\in \K}$ and we define the representation $G_{\qR_3}$ substituting $\K^n$ with $\K(x)^n$ (for any $n$) and $\lambda$ with $x$:
% https://q.uiver.app/#q=WzAsOCxbMCwxLCJcXEsoeCkiXSxbMSwxLCJcXEsoeCleMiJdLFsyLDEsIlxcSyh4KV4zIl0sWzMsMiwiXFxLKHgpXjIiXSxbMywwLCJcXEsoeCleMiJdLFs0LDEsIlxcSyh4KV4zIl0sWzUsMSwiXFxLKHgpXjIiXSxbNiwxLCJcXEsoeCkiXSxbMCwxLCJcXHBoaV97MSwxfSJdLFsxLDIsIlxccGhpX3syLDF9Il0sWzMsMiwiXFxwc2lfezIsMX0iXSxbMiw0LCJcXGJldGFfeCIsMCx7ImNvbG91ciI6WzM1OSwxMDAsNjBdfSxbMzU5LDEwMCw2MCwxXV0sWzMsNSwiXFxwaGlfezIsMX0iLDJdLFs1LDQsIlxcYWxwaGEiLDJdLFs2LDUsIlxccHNpX3syLDF9IiwyXSxbNyw2LCJcXHBzaV97MSwxfSIsMl0sWzMsNCwiIiwxLHsic3R5bGUiOnsiYm9keSI6eyJuYW1lIjoiZGFzaGVkIn0sImhlYWQiOnsibmFtZSI6Im5vbmUifX19XV0=
\[\begin{tikzcd}[ampersand replacement=\&,sep=small]
	\&\&\& {\K(x)^2} \&\&\& \\
	{\K(x)} \& {\K(x)^2} \& {\K(x)^3} \&\& {\K(x)^3} \& {\K(x)^2} \& {\K(x)} \\
	\&\&\& {\K(x)^2}
	\arrow["{\phi_{1,1}}", from=2-1, to=2-2]
	\arrow["{\phi_{2,1}}", from=2-2, to=2-3]
	\arrow["{\beta_x}", color={rgb,255:red,255;green,51;blue,54}, from=2-3, to=1-4]
	\arrow["\alpha"', from=2-5, to=1-4]
	\arrow["{\psi_{2,1}}"', from=2-6, to=2-5]
	\arrow["{\psi_{1,1}}"', from=2-7, to=2-6]
	\arrow[dashed, no head, from=3-4, to=1-4]
	\arrow["{\psi_{2,1}}", from=3-4, to=2-3]
	\arrow["{\phi_{2,1}}"', from=3-4, to=2-5]
\end{tikzcd}\]
with $\beta_x$ essentially being the previos map $\beta_\lambda$ in which we substitute the element $\lambda\in \K$ with the variable $x$.

\begin{lemma}
	$G_{\qR_3}$ is a generic brick over $\incidencealgebra{\qR_3}$.
\end{lemma}
\begin{proof}
	By \cite[Ch.III, Corollary 3.6]{sim:vol1}, we have that 
	\[\ell_{\incidencealgebra{\qR_3}}(G_{\qR_3}) = \dim_\K(G_{\qR_3}) = \sum_{i\in (\HasseQuiver{\qR_3})_0} \dim_\K((G_{\qR_3})_i) = \infty\]
	So $G_{\qR_3}$ has infinite length. Moreover, following the same procedure as in Lemma \ref{lemma:brickcontR3}, we see that all the components of an endomorphism $f\colon G_{\qR_3}\to G_{\qR_3}$ depends only on the map $f_1\colon \K(x) \to \K(x)$ that is $\K$-linear by definition and satisfies $xf_1(1)=f_1(x)$ (as done previously, substituting $\lambda$ with $x$). We also have that
	\[f_1(1)=f_1(xx^{-1})=xf_1(x^{-1}) \implies x^{-1}f_1(1)=f_1(x^{-1}) \]
	Therefore $f_1$ is completely determined by the element $f_1(1)\in \K(x)$, thus $\End(G_{\qR_3})\cong \K(x)$. Using again \cite[Ch.III, Corollary 3.6]{sim:vol1}
	\[\ell_{\K(x)\op}(G_{\qR_3}) = \dim_{\K(x)}(G_{\qR_3}) = \sum_{i\in (\HasseQuiver{\qR_3})_0} \dim_{\K(x)}((G_{\qR_3})_i) < \infty\]
	So $G_{\qR_3}$ has finite endolength too. By Remark \ref{rem:genmodulek(x)} we conclude that $G_{\qR_3}$ is a generic brick.
\end{proof}